\begin{center}
    {\LARGE \textbf{2018分析化学}} \\[2ex]
    {\large 时量180分钟 \quad 满分90分}
\end{center}

\vspace{0.5cm}
\noindent\textbf{注意事项:}
\begin{enumerate}[label=\arabic*., parsep=0pt, itemsep=0pt]
    \item 答题前,考生先将自己的姓名、学号、考场号、座位号填写在答题卡上,将条形码准确粘贴在考生信息条形码粘贴区。
    \item 考试结束后,将本试卷与答题纸一并收回。
    \item 回答各个问题时,将答案写在答题纸上,写在本试卷上无效。
\end{enumerate}

\vspace{0.5cm}
\noindent\textbf{一、简答题(42 pts.)}

\begin{enumerate}[label=\arabic*., itemsep=1.5ex]
    \item 请阐述银量法中吸附指示剂的变色原理，测定氯离子时为何不能用曙红为指示剂，若用曙红为指示剂对分析结果有何影响？
    \item 分析某混合碱，称取$0.3419\mathrm{g}$试样，溶于水后以酚酞为指示剂时，用$0.1000\mathrm{mol/L}\ \ce{HCl}$标准溶液滴定至终点，用去$\ce{HCl}$标准溶液$23.10\mathrm{mL}$。再加甲基橙指示剂继续滴定至终点，又消耗同浓度$\ce{HCl}$溶液$26.81\mathrm{mL}$，试问试样由何组分组成，各组分的含量为多少？（已知$M_{\ce{NaOH}}=40.01\mathrm{g\cdot mol^{-1}}$，$M_{\ce{Na2CO3}}=105.99\mathrm{g\cdot mol^{-1}}$，$M_{\ce{NaHCO3}}=84.01\mathrm{g\cdot mol^{-1}}$）
    \item 碘量法测定铜的实验中为何在接近终点时加入$\ce{KSCN}$？若不加对分析结果有何影响？
    \item 用$0.0200\mathrm{mol/L}$的$\ce{KMnO4}$溶液滴定$20.00\mathrm{mL}\ 0.1000\mathrm{mol/L}\ \ce{Fe^{2+}}$溶液。计算加入$15.00\mathrm{mL}$、$20.00\mathrm{mL}$、$22.00\mathrm{mL}\ \ce{KMnO4}$溶液时的电极电位。(已知$E^\ominus_{\ce{Fe^{3+}/Fe^{2+}}}=0.68\mathrm{V}$，$E^\ominus_{\ce{MnO4^-/Mn^{2+}}}=1.51\mathrm{V}$，$[\ce{H^+}]=1.0\mathrm{mol/L}$)
    \item 在下列物质中：
    \begin{itemize}
        \item $\ce{NH4Cl}$（$\mathrm{p}K_b(\ce{NH3})=4.74$）
        \item $\ce{C6H5OH}$（$\mathrm{p}K_a(\ce{C6H5OH})=9.96$）
        \item $\ce{Na2CO3}$（$\ce{H2CO3}$的$\mathrm{p}K_{a1}=6.38$，$\mathrm{p}K_{a2}=10.25$）
        \item $\ce{NaAc}$（$\mathrm{p}K_a(\ce{HAc})=4.74$）
        \item $\ce{HCOOH}$（$\mathrm{p}K_a(\ce{HCOOH})=3.74$）
    \end{itemize}
    请分别写出能用强碱标准溶液直接滴定的物质和能用强酸标准溶液直接滴定的物质。
    \item 写出下列溶液的质子条件式
    \begin{enumerate}[label=(\arabic*), noitemsep, topsep=0pt, partopsep=0pt]
        \item $0.1\ \mathrm{mol/L}$的$\ce{NH4H2PO4}$溶液
        \item $0.1\ \mathrm{mol/L}$的$\ce{H2SO4}$溶液
    \end{enumerate}
    \item 今以某弱酸$\ce{HA}$及其共轭碱$\ce{A^-}$配制缓冲溶液，已知其中$[\ce{HA}]=0.25\mathrm{mol/L}$。往上述$100\mathrm{mL}$缓冲溶液中加入$5.0\mathrm{mmol}$固体$\ce{NaOH}$后(忽略其体积的变化)，溶液的$\mathrm{pH}=4.60$。计算原缓冲溶液的$\mathrm{pH}$(设$\ce{HA}$的$\mathrm{p}K_a=4.30$)。
\end{enumerate}

\vspace{0.5cm}
\noindent\textbf{二、计算题(48 pts.)}

\begin{enumerate}[label=\arabic*., itemsep=1.5ex]
    \item (10分) 今有某弱酸$\ce{HA}$，其$\mathrm{p}K_a=5.5$。若以$0.2000\mathrm{mol/L}\ \ce{NaOH}$溶液滴定$20.00\mathrm{mL}\ 0.2000\mathrm{mol/L}$的该弱酸溶液，试计算滴定曲线的突跃范围及化学计量点的$\mathrm{pH}$值。如用酚酞作指示剂，滴定至$\mathrm{pH}=9.0$时为终点，试计算终点误差。
    \item (10分) 以$0.02000\mathrm{mol/L}\ \mathrm{EDTA}$滴定浓度均为$0.0200\mathrm{mol/L}\ \ce{Pb^{2+}}$，$\ce{Ca^{2+}}$混合液中的$\ce{Pb^{2+}}$，溶液$\mathrm{pH}$为$5.0$。计算：
    \begin{enumerate}[label=(\arabic*), noitemsep, topsep=0pt, partopsep=0pt]
        \item 化学计量点时的$\lg K'(\ce{PbY})$和$[\ce{Pb^{2+}}]$值
        \item 若以二甲酚橙为指示剂，终点误差多大？已知$\mathrm{pH}=5.0$时$\lg\alpha_{Y(\mathrm{H})}=6.6$，$\mathrm{pPb_{终}}=7.0$(二甲酚橙)，$\lg K(\ce{PbY})=18.0$，$\lg K(\ce{CaY})=10.7$。
    \end{enumerate}
    \item (10分) 当溶液中$[\ce{Cl^-}]=0.010\mathrm{mol/L}$，$[\ce{CrO4^{2-}}]=0.10\mathrm{mol/L}$时，逐滴加入$\ce{AgNO3}$溶液，首先生成哪一种沉淀？当第二种离子开始沉淀时，第一种未被沉淀的离子浓度是多少？(已知$K_{\mathrm{sp,AgCl}}=1.8\times10^{-10}$，$K_{\mathrm{sp,Ag2CrO4}}=1.1\times10^{-12}$)
    \item (8分) 溶解含铜($M_r=63.55\mathrm{g\cdot mol^{-1}}$)试样$0.1000\mathrm{g}$，定容$50\mathrm{mL}$。分取$25.00\mathrm{mL}$，用有机螯合萃取溶剂$25.00\mathrm{mL}$将铜萃入有机相，分配比是$58$。在适当条件，用$1\mathrm{cm}$比色皿测得有机相吸光度为$0.635$，摩尔吸光系数为$5.72\times10^3\mathrm{L\cdot mol^{-1}\cdot cm^{-1}}$，问试样中铜的含量是多少？如果仅用有机相测得结果计算试样中铜含量，相对误差有多大？
    \item (10分) 用碘量法测定某铜合金中铜的质量分数$\omega(\ce{Cu})/\%$，$6$次测定结果如下：$60.60$，$60.64$，$60.58$，$60.65$，$60.57$和$60.32$。(1) 用格鲁布斯法检验有无应舍弃的异常值（显著水平$0.05$）；(2) 估计铜的质量分数范围($P=95\%$)；(3) 如果铜的质量分数标准值为$60.58\%$，试问测定有无系统误差(显著水平$0.05$)？
    \\
    已知：
    \begin{center}
    \begin{tabular}{l|ccc}
    $n$ & $4$ & $5$ & $6$ \\
    \hline
    $T_{0.05}$ & $1.463$ & $1.672$ & $1.832$ \\
    \end{tabular}
    \begin{tabular}{l|ccc}
    $f$ & $4$ & $5$ & $6$ \\
    \hline
    $t_{0.05}$ & $2.776$ & $2.571$ & $2.447$ \\
    \end{tabular}
    \end{center}
\end{enumerate}